%% Phase 1: Jamaratv9 Parameter Tables for §5.3
%% Extracted from MFGraphs/Examples/ExamplesData.wl lines 302-318
%% Parameters: I1->100, I2->50, U1->0, U2->0, U3->0 (from BenchmarkHelpers.wls $DefaultParams)

%%% TABLE 1: Network Topology Summary
\begin{table}[h!]
\caption{Jamaratv9 network topology}
\label{tab:jamaratv9-topology}
\centering
\begin{tabular}{ll}
\toprule
\textbf{Property} & \textbf{Value} \\
\midrule
Vertices & $|\mathcal{V}| = 9$ \\
Directed edges & $|\mathcal{E}| = 11$ \\
Entrance vertices & 2 (vertices 1, 2) \\
Exit vertices & 3 (vertices 7, 8, 9) \\
Switching costs & None \\
\bottomrule
\end{tabular}
\end{table}

%%% TABLE 2: Edge List
\begin{table}[h!]
\caption{Jamaratv9 directed edges}
\label{tab:jamaratv9-edges}
\centering
\begin{tabular}{cccc}
\toprule
\textbf{Source} & \textbf{Target} & \textbf{Source} & \textbf{Target} \\
\midrule
$1$ & $2$ & $5$ & $6$ \\
$1$ & $3$ & $5$ & $7$ \\
$2$ & $4$ & $6$ & $8$ \\
$3$ & $4$ & $7$ & $9$ \\
$3$ & $5$ & $8$ & $9$ \\
$4$ & $6$ &  &  \\
\bottomrule
\end{tabular}
\end{table}

%%% TABLE 3: Entrance Flow and Exit Terminal Costs
\begin{table}[h!]
\caption{Jamaratv9 entrance flows and exit terminal costs}
\label{tab:jamaratv9-flow-costs}
\centering
\begin{tabular}{lcc|lcc}
\toprule
\multicolumn{3}{c}{\textbf{Entrance flows}} & \multicolumn{3}{c}{\textbf{Exit terminal costs}} \\
\textbf{Vertex} & \textbf{Symbol} & \textbf{Value} & \textbf{Vertex} & \textbf{Symbol} & \textbf{Value} \\
\midrule
$1$ & $I_1$ & $100$ & $7$ & $U_1$ & $0$ \\
$2$ & $I_2$ & $50$ & $8$ & $U_2$ & $0$ \\
 &  &  & $9$ & $U_3$ & $0$ \\
\bottomrule
\end{tabular}
\end{table}

%%% TABLE 4: Problem Specification Summary
\begin{table}[h!]
\caption{Jamaratv9 alternatives and pruning target}
\label{tab:jamaratv9-alternatives}
\centering
\begin{tabular}{lp{4cm}}
\toprule
\textbf{Quantity} & \textbf{Description} \\
\midrule
Equilibrium sign triples & $q$ (measured during preprocessing)* \\
Total sign alternatives & $2^q$ \\
Expected count & $2^{95}$ (to be measured) \\
Solver goal & Identify surviving leaves $M$ among $2^q$ branches \\
\bottomrule
\multicolumn{2}{p{11cm}}{\footnotesize *The value $q$ is the number of ``uncertain'' switching triples $(i,j,k)$ where the optimality condition is non-strict (either direction feasible). See \S\ref{sec:methodology} for details.} \\
\end{tabular}
\end{table}

%%% NOTES FOR §5.3
% Insert these tables into §5.3 before the pruning statistics discussion.
% Cross-check edge directions and labels with figures/Jama.pdf
% After running the solver in Phase 4, measure the actual q value and update "Expected count" in Table 4
% If q ≠ 95, update the abstract and conclusion accordingly
